Celkové ztráty při přemagnetování jsou $P=P_h+P_v+P_z$:
\begin{itemize}
\item \textbf{Hysterezní $P_h$:} souvisejí se změnami domén; $P_h=VS_hf$, tedy při stejné smyčce rostou \textbf{lineárně s frekvencí}. $V$ je objem a $S_h$ plocha statické smyčky.
\item \textbf{Vířivými proudy $P_v$:} indukované proudy ohřívají elektricky vodivé jádro. Při stejné geometrii, rezistivitě a amplitudě indukce platí \textbf{$P_v\propto f^2$}.
\item \textbf{Magnetickým zpožděním $P_z$:} souvisejí s pohyby elektronů v mřížce při změnách pole; významné jsou u feritů a v malých polích u měkkých materiálů. Přednáška pro ně konkrétní kmitočtovou závislost neuvádí.
\end{itemize}
